\section{Hardware Architecture}
\label{sec:arch}

% The core is scalar + vector as one pipeline. XDAI / GEMM after.
% Pipeline figure shows the shared frontend. Die lives in Impl
% (first cite). No clock/area here.

\paragraph{RISC-V Core.}
The core is an in-order, single-issue, four-stage pipeline, divided into a shared frontend and a scalar and vector backends.
The frontend stages are Instruction Fetch Unit ({IFU}), Decode~0 ({IDU0}),
Decode~1 ({IDU1}), while the backends are the scalar and vector Execute ({EXU}) as depicted in Figure~\ref{fig:core-pipe}.

\begin{figure}[t]
  \centering
  % Single-column core pipeline. Theme match pyvedas_flow.
% XRF in frontend (read at IDU1). Writeback via a small WB hub
% under the backends, then one rail up into XRF.
\definecolor{tvink}{HTML}{111111}
\definecolor{tvmute}{HTML}{333333}
\definecolor{tvline}{HTML}{555555}

\pgfdeclarelayer{backboxes}
\pgfdeclarelayer{backarrows}
\pgfsetlayers{backboxes,backarrows,main}

\begin{tikzpicture}[
  >=Stealth,
  x=1mm, y=1mm,
  font=\sffamily\tiny,
  box/.style={
    draw=tvink,
    line width=0.55pt,
    rounded corners=1.2pt,
    fill=white,
    align=center,
    inner xsep=2.0pt,
    inner ysep=1.7pt,
    font=\sffamily\tiny\bfseries,
  },
  unit/.style={
    box,
    font=\sffamily\tiny,
    inner xsep=1.8pt,
    inner ysep=1.6pt,
    minimum width=9.2mm,
  },
  arr/.style={
    ->,
    draw=tvline,
    line width=0.5pt,
    shorten >=0.5pt,
    shorten <=0.5pt,
  },
]
\pgfmathsetmacro{\W}{\columnwidth/1mm - 0.8}
\pgfmathsetmacro{\exuB}{10.0}
\pgfmathsetmacro{\exuT}{33.2}
\pgfmathsetmacro{\exuM}{0.5*(\exuB+\exuT)}

\draw[tvink, line width=0.55pt, rounded corners=2pt]
  (0, 0.4) rectangle (\W, 54.2);

\node[font=\sffamily\scriptsize\bfseries, text=tvink, anchor=west]
  at (1.6, 51.6) {Core pipeline};
\draw[tvink, line width=0.4pt]
  (1.6, 49.6) -- (\W-1.6, 49.6);

\node[font=\sffamily\tiny\bfseries, text=tvink, anchor=west]
  at (1.6, 46.8) {Shared frontend};

\node[box] (ifu)  at (0.14*\W, 40.8) {IFU};
\node[box] (idu0) at (0.36*\W, 40.8) {IDU0};
\node[box] (idu1) at (0.58*\W, 40.8) {IDU1};
\node[box, line width=0.85pt, align=center,
  inner ysep=1.4pt] (xrf) at (0.84*\W, 40.8)
  {XRF\\[-0.15em]{\normalfont\sffamily\tiny\color{tvmute}integer RF}};

\begin{pgfonlayer}{backboxes}
  \draw[tvink, densely dotted, line width=0.55pt, rounded corners=2pt]
    (1.4, \exuB) rectangle (\W-4.2, \exuT);
\end{pgfonlayer}
\node[font=\sffamily\tiny\bfseries, text=tvink, fill=white,
  inner xsep=1.6pt, inner ysep=0.2pt]
  at (\W-6.2, \exuT) {EXU};

\node[font=\sffamily\tiny\bfseries, text=tvink, anchor=west]
  (lab_s) at (2.2, \exuT-2.6) {Scalar};
\node[font=\sffamily\tiny\bfseries, text=tvink, anchor=west]
  (lab_v) at (0.52*\W, \exuT-2.6) {Vector};
\draw[tvink, line width=0.4pt] (2.2, \exuT-4.0) -- (0.46*\W, \exuT-4.0);
\draw[tvink, line width=0.4pt] (0.52*\W, \exuT-4.0) -- (\W-4.8, \exuT-4.0);

\node[unit] (alu) at (0.16*\W, \exuM+4.6) {ALU};
\node[unit] (mul) at (0.34*\W, \exuM+4.6) {MUL};
\node[unit] (div) at (0.16*\W, \exuM-4.6) {DIV};
\node[unit] (lsu) at (0.34*\W, \exuM-4.6) {LSU};

\node[unit] (valu) at (0.62*\W, \exuM+4.6) {VALU};
\node[unit] (vlsu) at (0.80*\W, \exuM+4.6) {VLSU};
\node[unit] (vrf)  at (0.62*\W, \exuM-4.6) {VRF};
\node[unit] (vcsr) at (0.80*\W, \exuM-4.6) {vset};

\node[box, font=\sffamily\tiny\bfseries, inner xsep=3pt,
  inner ysep=1.3pt, minimum width=12mm] (wb) at (0.50*\W, 3.4) {WB};

\begin{pgfonlayer}{backarrows}
  \draw[arr] (ifu) -- (idu0);
  \draw[arr] (idu0) -- (idu1);
  \draw[arr] (xrf.west) -- (idu1.east);
  \draw[arr] (idu1.south) -- ++(0,-1.4) -| (0.25*\W, \exuT-4.0);
  \draw[arr] (idu1.south) -- ++(0,-1.4) -| (0.71*\W, \exuT-4.0);
  \draw[tvline, line width=0.5pt]
    (0.25*\W, \exuB-0.2) -- (0.25*\W, \exuB-1.4);
  \draw[tvline, line width=0.5pt]
    (0.71*\W, \exuB-0.2) -- (0.71*\W, \exuB-1.4);
  \draw[tvline, line width=0.5pt]
    (0.25*\W, \exuB-1.4) -- (0.71*\W, \exuB-1.4);
  \draw[arr] (0.50*\W, \exuB-1.4) -- (wb.north);
  \draw[arr] (wb.east) -- (\W-2.4, 3.4) -- (\W-2.4, 40.8) -- (xrf.east);
\end{pgfonlayer}

\end{tikzpicture}

  \caption{Core pipeline: shared {IFU}/{IDU0}/{IDU1} with the
    integer scalar register file ({XRF}); scalar and vector {EXU};
    both write {XRF}, while the vector register file ({VRF}) stays in
    the vector datapath.}
  \label{fig:core-pipe}
\end{figure}

{IFU} maintains a program counter into the Instruction Closely Coupled
Memory ({ICCM}).
{IDU0} runs the {RV32IM} and {Zve32x} decoders in parallel on the same
instruction and dispatches to either the scalar or vector backend.
{IDU1} reads the integer scalar register file ({XRF}), checks the
scoreboard, and
issues to {EXU}, or stalls the frontend of the machine.

The scalar datapath is composed by an integer Arithmetic-Logic Unit ({ALU}), a multiply unit, 
a divide unit, and the scalar load/store unit ({LSU}).
The {LSU} is the bridge between the core and the XDAI, connecting it to the available DAs
and the memory subsystem.

The vector datapath implements {Zve32x} specifications, with 512-bit vectors
({VLEN}), 32-bit elements and unit-stride loads and stores.
To keep area utilization low, the vector {ALU} and memory datapath are
only 128~bits wide ({DLEN}) rather than a full 512-bit slice, so a
vector op walks the register in four beats.

\paragraph{eXtensible Datapath Accelerator Interconnect.} % XDAI
The {XDAI} is the Advanced eXtensible Interface ({AXI}) fabric the
{LSU} uses to reach memories, peripherals, and {MMIO}
{DA}s. A software-defined range decoder matches each address against the
device map and steers the memory request to the target device.

This unified bus allows a {DA} to expose its Control and Status Registers ({CSRs})
to the RISC-V core. Once triggered by the core, a {DA} can take control of the
{XDAI} and access any memory-mapped device, such as Data Closely Coupled Memory ({DCCM}) and other {SRAM} resources,
where weights might be stored, for example, through a Direct Memory Access ({DMA}) engine.

The {DA} gives control back to the core by setting the {Done} {CSR}.

The memory map is software-defined.
A device-tree-style YAML file lists each {DA}'s address window.
Hardware uses it to build the {SoC} decode; the same file emits the
{C} macros the {XCL} includes when it programs a {DA}.


\paragraph{Datapath Accelerators.}
DAs are memory-mapped, domain-specific hardware blocks used by the core to offload a
particular functionality to specialized hardware.
In this paper, the only {DA} on the {XDAI} is an \texttt{int8} $8{\times}8$ integer
General Matrix Multiply ({GEMM}) engine, performing $A \times B = C$. 

The array is output-stationary $8{\times}8$ {PE}s with $K{=}32$.
Software writes the three bases---including $C$---and $M$, $N$, $K$,
then toggles the {START} {CSR}.

{GEMM}'s {DMA} reads $A$ and $B$ from {DCCM} and writes $C$ to the
base the core programmed, then sets the {Done} {CSR}.
Figure~\ref{fig:soc-block} is the {SoC}: the core, the {XDAI} with that
{GEMM}, and the ordinary peripherals---Universal Asynchronous Receiver--Transmitter ({UART}), end-of-test ({EOT}),
and the on-chip memories.

\begin{figure}[H]
  \centering
  % Single-column SoC block: core → XDAI → devices / memories.
% Theme match pyvedas_flow / core_pipeline.
\definecolor{tvink}{HTML}{111111}
\definecolor{tvmute}{HTML}{333333}
\definecolor{tvline}{HTML}{555555}
\definecolor{tvshade}{HTML}{EEEEEE}

\pgfdeclarelayer{backarrows}
\pgfsetlayers{backarrows,main}

\begin{tikzpicture}[
  >=Stealth,
  x=1mm, y=1mm,
  font=\sffamily\scriptsize,
  box/.style={
    draw=tvink,
    line width=0.6pt,
    rounded corners=1.5pt,
    fill=white,
    align=center,
    inner xsep=2.5pt,
    inner ysep=2.2pt,
    font=\sffamily\scriptsize,
  },
  hub/.style={
    box,
    line width=1.15pt,
    fill=tvshade,
    font=\sffamily\footnotesize\bfseries,
  },
  arr/.style={
    ->,
    draw=tvline,
    line width=0.55pt,
    shorten >=0.8pt,
    shorten <=0.8pt,
  },
]
\pgfmathsetmacro{\W}{\columnwidth/1mm - 0.8}

\draw[tvink, line width=0.55pt, rounded corners=2pt]
  (0, 0.6) rectangle (\W, 58.5);

\node[font=\sffamily\footnotesize\bfseries, text=tvink, anchor=west]
  at (1.6, 55.8) {SoC};
\draw[tvink, line width=0.4pt]
  (1.6, 53.8) -- (\W-1.6, 53.8);

\node[box, minimum width=0.18*\W] (iccm) at (0.14*\W, 48.2) {ICCM};
\node[hub, minimum width=0.48*\W, align=center,
  inner ysep=2.2pt] (core) at (0.5*\W, 48.2)
  {Core\\[0.12em]
   {\normalfont\sffamily\scriptsize\color{tvmute}scalar + vector}};

\node[hub, minimum width=0.72*\W, inner ysep=2.0pt] (bus)
  at (0.5*\W, 39.6)
  {XDAI};

\node[font=\sffamily\scriptsize\bfseries, text=tvink, anchor=west]
  at (1.6, 33.4) {Devices};
\node[box, minimum width=0.24*\W] (uart) at (0.22*\W, 28.4)
  {UART\\[-0.1em]{\normalfont\sffamily\tiny\color{tvmute}0x00200000}};
\node[box, line width=0.9pt, minimum width=0.24*\W] (gemm) at (0.50*\W, 28.4)
  {{GEMM}\\[-0.1em]{\normalfont\sffamily\tiny\color{tvmute}0x00300000}};
\node[box, minimum width=0.24*\W] (eot) at (0.78*\W, 28.4)
  {EOT\\[-0.1em]{\normalfont\sffamily\tiny\color{tvmute}0x10000000}};

\node[font=\sffamily\scriptsize\bfseries, text=tvink, anchor=west]
  at (1.6, 21.6) {Memories};
\node[box, minimum width=0.24*\W] (dccm) at (0.50*\W, 11.0) {DCCM};
\node[box, minimum width=0.24*\W] (sram) at (0.78*\W, 11.0)
  {SRAM stub\\[-0.1em]{\normalfont\sffamily\tiny\color{tvmute}0x40000000}};

\begin{pgfonlayer}{backarrows}
  \draw[arr] (core.west) -- (iccm.east);
  \draw[arr] (core.south) -- (bus.north);

  \draw[tvline, line width=0.5pt]
    (bus.south) -- ++(0,-2.0) coordinate (devrail);
  \draw[tvline, line width=0.5pt]
    (uart.north |- devrail) -- (eot.north |- devrail);
  \draw[arr] (uart.north |- devrail) -- (uart.north);
  \draw[arr] (gemm.north |- devrail) -- (gemm.north);
  \draw[arr] (eot.north |- devrail) -- (eot.north);

  % Memories leave the XDAI on the east side, then drop.
  \draw[tvline, line width=0.55pt]
    (bus.east) -- (0.92*\W, 0 |- bus.east) -- (0.92*\W, 19.8);
  \draw[tvline, line width=0.55pt]
    (0.50*\W, 19.8) -- (0.92*\W, 19.8);
  \draw[arr] (0.50*\W, 19.8) -- (dccm.north);
  \draw[arr] (0.78*\W, 19.8) -- (sram.north);
\end{pgfonlayer}

\end{tikzpicture}

  \caption{{SoC} block: core, {MMIO}/{XDAI}, {GEMM} as the attached
    {DA}, {UART}, {EOT}, and memories.}
  \label{fig:soc-block}
\end{figure}


% Memory map: on-chip SRAM stub @ 0x40000000 stands in for DRAM (FPGA:
% URAM/BRAM in axi4_dram_fpga --- say SRAM in the PDF, not UltraRAM).

\paragraph{Plugging the accelerator into \pyv{}.}
Once the {DA} is on the {XDAI}, \pyv{} still has to see it:
an \texttt{ops.yaml} row, a codegen handler, an {XCL} body, and the
hardware {YAML} for the address window.

The first step is to write a codegen handler for the accelerator.
That handler is the compiler's picture of the {CSR} map: a base pointer
for $A$, one for $B$, one for $C$, the three shapes $M$, $N$, and $K$,
and a {START} bit.
Those registers live in the {GEMM}'s window on the software-defined
memory map, so that the core's LSU can directly access those registers.

The next step is to review the primitives inside {XCL} and ask a simple question: which primitives
could actually benefit from this {DA}?
Matrix multiply is the obvious one.
Convolution is the useful one, because a convolution is a matrix
multiply after an \texttt{im2col}.
Those {XCL} bodies used to be native {C} loops, translated into RISC-V code by GCC.
To offload matrix multiplication to {GEMM}, I just substitute the inner multiply loop of the convolution
with a call that programs the {GEMM} and lets its {DMA} engine walk $A$ and $B$.

The last step is compile-time tiling.
The engine is $8{\times}8$ with $K{=}32$.
\pyv{} cuts a large multiply when the packed $A$/$B$ panels would
not fit on-chip; the generated {C} walks those tiles, one {START}
at a time.
Figure~\ref{fig:pyvedas-gemm-plug} is that whole path for this {GEMM}.

\begin{figure}[H]
  \centering
  % How PyVedas learns the GEMM DA. Column-width, print theme.
\definecolor{tvink}{HTML}{111111}
\definecolor{tvmute}{HTML}{333333}
\definecolor{tvline}{HTML}{555555}
\definecolor{tvshade}{HTML}{EEEEEE}

\pgfdeclarelayer{backarrows}
\pgfsetlayers{backarrows,main}

\begin{tikzpicture}[
  >=Stealth,
  x=1mm, y=1mm,
  font=\sffamily\tiny,
  box/.style={
    draw=tvink,
    line width=0.55pt,
    rounded corners=1.2pt,
    fill=white,
    align=left,
    inner xsep=2.4pt,
    inner ysep=2.0pt,
    font=\ttfamily\fontsize{5.2}{6.4}\selectfont,
  },
  stepnum/.style={
    draw=tvink,
    line width=0.85pt,
    rounded corners=1.2pt,
    fill=tvshade,
    font=\sffamily\tiny\bfseries,
    inner xsep=2.2pt,
    inner ysep=1.2pt,
  },
  arr/.style={
    ->,
    draw=tvline,
    line width=0.5pt,
    shorten >=0.5pt,
    shorten <=0.5pt,
  },
]
\pgfmathsetmacro{\W}{\columnwidth/1mm - 0.8}

\draw[tvink, line width=0.55pt, rounded corners=2pt]
  (0, -1.0) rectangle (\W, 72.0);

\node[font=\sffamily\scriptsize\bfseries, text=tvink, anchor=west]
  at (1.6, 69.0) {Plugging the {GEMM} into \pyv{}};
\draw[tvink, line width=0.4pt]
  (1.6, 67.0) -- (\W-1.6, 67.0);

\node[stepnum] (s1) at (0.12*\W, 61.5) {1};
\node[box, text width=0.72*\columnwidth, anchor=west]
  (c1) at (0.22*\W, 61.5) {%
  \textcolor{tvmute}{\# hw/soc/default.yaml}\\
  - name: gemm\\
  \quad role: accelerator\\
  \quad module: gemm\_top%
};

\node[stepnum] (s2) at (0.12*\W, 46.5) {2};
\node[box, text width=0.72*\columnwidth, anchor=west]
  (c2) at (0.22*\W, 46.5) {%
  \textcolor{tvmute}{\# ops.yaml}\\
  pyvedas.gemm\_mmio.default:\\
  \quad codegen: gemm\_mmio\\
  \quad symbol: pyvedas\_gemm\_job\\
  \quad sources: gemm\_mmio.c%
};

\node[stepnum] (s3) at (0.12*\W, 29.5) {3};
\node[box, text width=0.72*\columnwidth, anchor=west]
  (c3) at (0.22*\W, 29.5) {%
  \textcolor{tvmute}{\# codegen\_handlers.py}\\
  emit\_gemm\_mmio(...)\\
  \quad $\rightarrow$ calls in generated C%
};

\node[stepnum] (s4) at (0.12*\W, 12.5) {4};
\node[box, text width=0.72*\columnwidth, anchor=west]
  (c4) at (0.22*\W, 12.5) {%
  \textcolor{tvmute}{\# gemm\_mmio.c (XCL)}\\
  write BASE\_A/B/C, M, N, K\\
  \quad store START on MMIO\_GEMM\_ADDR%
};

\begin{pgfonlayer}{backarrows}
  \draw[arr] (c1.south) -- (c2.north);
  \draw[arr] (c2.south) -- (c3.north);
  \draw[arr] (c3.south) -- (c4.north);
\end{pgfonlayer}

\end{tikzpicture}

  \caption{Plugging the {GEMM} into \pyv{}: a codegen handler for the
    {CSR} map, an {XCL} body that calls the {DA}, and a tiling pass
    so packed panels fit on-chip.}
  \label{fig:pyvedas-gemm-plug}
\end{figure}
